\documentclass[11pt]{article}

\usepackage[T1]{fontenc}
\usepackage[utf8]{inputenc}
\usepackage{amsmath,amssymb,amsthm,mathtools}
\usepackage{booktabs}
\usepackage[margin=1in]{geometry}
\usepackage{microtype}
\usepackage{hyperref}
\usepackage{xcolor}

\hypersetup{
  colorlinks=true,
  linkcolor=blue!45!black,
  citecolor=blue!45!black,
  urlcolor=blue!45!black,
  pdftitle={Exact reproduction and KS-atlas transfer of a Morales-Ramis obstruction near the parabolic Lagrange solution},
  pdfauthor={Jacob Sullivan}
}

\newtheorem{theorem}{Theorem}
\newtheorem{proposition}[theorem]{Proposition}
\newtheorem{lemma}[theorem]{Lemma}
\newtheorem{corollary}[theorem]{Corollary}
\theoremstyle{definition}
\newtheorem{definition}[theorem]{Definition}
\theoremstyle{remark}
\newtheorem{remark}[theorem]{Remark}

\newcommand{\C}{\mathbb{C}}
\newcommand{\R}{\mathbb{R}}
\newcommand{\T}{\mathbb{T}}
\newcommand{\XiMap}{\Xi}
\newcommand{\Id}{\mathrm{I}}
\newcommand{\Span}{\operatorname{span}}
\newcommand{\rank}{\operatorname{rank}}
\newcommand{\Gal}{\operatorname{Gal}}

\title{\textbf{Exact reproduction and KS-atlas transfer of a
Morales--Ramis obstruction near the parabolic Lagrange solution}}
\author{Jacob Sullivan\\
\small Affiliation and ORCID to be confirmed before submission}
\date{24 July 2026}

\begin{document}
\maketitle

\begin{abstract}
We give an exact, machine-qualified reproduction of the planar
Morales--Ramis obstruction near the parabolic Lagrange phase curve for every
finite positive mass triple.  Starting from the declared
translation- and rotation-reduced Hamiltonian, an implementation disjoint
from the production certificate differentiates the Hamiltonian, verifies the
phase curve, constructs the full six-dimensional variational equation,
performs the Whittaker energy-shell reduction, and derives a four-dimensional
Fuchsian normal variational equation.  Its finite residue spectra and
infinity characteristic polynomial imply that the differential-Galois
identity component is non-Abelian throughout the positive-mass domain.  We
then transfer that local obstruction to the translation-reduced spatial
Hamiltonian by an explicit variational subquotient and to a qualified
two-chart Kustaanheimo--Stiefel (KS) atlas by exact cotangent reduction,
fixed-energy Sundman equivalence, and a shell-relative rank argument.
The resulting theorem concerns complete meromorphic Liouville integrability
locally near one fixed-nonzero-angular-momentum phase curve.  It does not
exclude one additional integral or make a global assertion.  The planar
theorem is a reproduction of published work, not an asserted original
discovery.  All decision-bearing formulas, source hashes, negative controls,
and qualification records are included in a one-command reproducibility
package.
\end{abstract}

\section{Scope and prior art}
\label{sec:scope}

Differential Galois theory supplies necessary conditions for the meromorphic
Liouville integrability of Hamiltonian systems.  In the form used here, the
Morales--Ramis theorem requires the identity component of the differential
Galois group of the variational equation, and hence of the relevant normal
variational equation, to be Abelian near a non-equilibrium phase curve
\cite{moralesramis2001}.  Tsygvintsev applied this framework to the
parabolic Lagrange solution of the planar three-body problem and derived the
Fuchsian normal system reproduced below \cite{tsygvintsev2002}.
Boucher and Weil gave an independent published application at the theorem
level \cite{boucherweil}.

The scientific contribution qualified in this paper has four deliberately
separated parts:
\begin{enumerate}
  \item an exact project-owned reconstruction of the planar normal
    variational equation from the reduced Hamiltonian;
  \item a source-equivalence audit against hash-pinned primary papers;
  \item an exact transfer of the obstruction to the spatial variational
    module; and
  \item an exact transfer to a previously qualified multi-chart KS
    regularization.
\end{enumerate}
The planar nonintegrability result itself is prior art.  We do not claim its
discovery, its first proof, or novelty of its mathematical conclusion.

The word ``integrability'' will always mean complete Liouville integrability
by meromorphic first integrals on a complex neighborhood unless another
meaning is explicitly stated.  For the translation-reduced spatial
Hamiltonian, which has six degrees of freedom, completeness requires six
functionally independent commuting integrals including the Hamiltonian.
The result is local to the neighborhood of a specified phase curve.  It is
not an assertion about every trajectory or every phase-space domain.

\subsection{Models and normalizations}

All masses are finite and strictly positive.  In the planar reproduction we
set the gravitational constant and the fixed angular-momentum scale to one,
and normalize the third mass to \(m_3=1\).  Positive common scaling recovers
every positive mass triple.  The spatial theorem is stated on the regular
collision-free portion of the translation-reduced phase space.  Binary
collision extensions enter only through restriction back to that open set;
triple collision is outside the qualified model.

\section{Reduced planar model}
\label{sec:planar-model}

Let
\[
  M_1=1+\frac{1}{m_1},\qquad
  M_2=1+\frac{1}{m_2},\qquad
  P=p_3q_2-p_2q_3-1.
\]
The declared reduced Hamiltonian is
\begin{align}
H={}&\frac{M_1}{2}\left(p_1^2+\frac{P^2}{q_1^2}\right)
  +\frac{M_2}{2}(p_2^2+p_3^2)+p_1p_2-\frac{p_3P}{q_1}
  \nonumber\\
 &-\frac{m_1}{q_1}
  -\frac{m_2}{\sqrt{q_2^2+q_3^2}}
  -\frac{m_1m_2}{\sqrt{(q_1-q_2)^2+q_3^2}} .
\label{eq:reduced-H}
\end{align}
The implementation reconstructs every derivative used below from
\eqref{eq:reduced-H}; it does not import a published variational matrix.

Put
\[
 S_1=m_1+m_2+1,\qquad
 S_2=m_1m_2+m_1+m_2,\qquad
 S_3=m_2+2.
\]
The equilateral parabolic phase curve is parametrized by
\begin{align}
 q_1&=q,&q_2&=\frac q2,&q_3&=\frac{\sqrt3\,q}{2},\nonumber\\
 p_1&=p,&p_2&=Ap+\frac Bq,&p_3&=Cp+\frac Dq,
\label{eq:lagrange-curve}
\end{align}
where
\begin{align*}
A&=\frac{m_2(1-m_1)}{m_1S_3},&
B&=-\frac{\sqrt3\,S_1m_2}{S_2S_3},\\
C&=\frac{\sqrt3\,m_2(m_1+1)}{m_1S_3},&
D&=-\frac{m_2(S_2+m_1m_2-1)}{S_2S_3}.
\end{align*}
On the zero-energy surface let \(w=pq\).  Then
\begin{equation}
 q(w)=-\frac{aw^2+bw+d}{c},\qquad p(w)=\frac{w}{q(w)},
\label{eq:q-of-w}
\end{equation}
with
\[
 a=\frac{2S_1S_2}{m_1^2S_3^2},\quad
 b=-\frac{2\sqrt3\,m_2S_1}{m_1S_3^2},\quad
 c=-S_2,\quad
 d=\frac{2S_1(m_2^2+m_2+1)}{S_3^2S_2}.
\]

\begin{proposition}[Exact phase-curve check]
\label{prop:phase}
For every positive \(m_1,m_2\), equations
\eqref{eq:lagrange-curve}--\eqref{eq:q-of-w} define a non-equilibrium
zero-energy integral phase curve of \eqref{eq:reduced-H} away from its
isolated singular parameter values.
\end{proposition}

\begin{proof}
Direct substitution gives \(H=0\).  The exact reference calculation also
verifies that all six components of \(J_6\nabla H\) equal the corresponding
components of \(dz/dw\) multiplied by one common rational expression
\(dw/dt\).  This stronger tangent identity rules out the weaker possibility
that the displayed parametrization merely lies on the energy shell.
\end{proof}

\section{Exact normal variational equation}
\label{sec:nve}

The full planar variational equation along the phase curve is
\[
 \frac{d\,\delta z}{dt}
   =J_6 H_{zz}(\Gamma(w))\,\delta z .
\]
The reference implementation forms this \(6\) by \(6\) matrix by exact
differentiation and verifies its Hamiltonian-matrix identity.

To remove the energy/flow pair, solve \(H=0\) locally for \(p_1\) and put
\(K=-p_1\).  For \(y=(q_2,q_3,p_2,p_3)\), implicit differentiation gives
\begin{equation}
 K_{ij}=\frac{H_{ij}}{H_p}
 -\frac{H_{ip}H_j}{H_p^2}
 -\frac{H_iH_{pj}}{H_p^2}
 +\frac{H_iH_{pp}H_j}{H_p^3},
\label{eq:whittaker-hessian}
\end{equation}
where \(p=p_1\).  After substituting the phase curve and changing the
independent variable from \(q_1\) to \(w\), the four-dimensional normal
system is
\begin{equation}
 A_{\mathrm{NVE}}
   =J_4\,\frac{dq}{dw}\,\operatorname{Hess}_y(K).
\label{eq:nve-matrix}
\end{equation}
Symmetry of the Hessian and the Hamiltonian identity for
\eqref{eq:nve-matrix} are exact qualification gates.

Define
\begin{align}
 L&=2S_2w-\sqrt3\,m_1m_2,\nonumber\\
 Z&=S_2^2w^2-\sqrt3\,m_1m_2S_2w
      +m_1^2(m_2^2+m_2+1),\nonumber\\
 C&=\operatorname{diag}(LZ,LZ,1,1).
\label{eq:gauge}
\end{align}
The gauge-transformed system is
\[
 A_F=C^{-1}A_{\mathrm{NVE}}C-C^{-1}\frac{dC}{dw}.
\]
Exact denominator valuations show that it is Fuchsian, with three distinct
finite singularities
\begin{align}
t_0&=\frac{\sqrt3\,m_1m_2}{2S_2},\nonumber\\
t_1&=\frac{m_1(\sqrt3\,m_2+iS_3)}{2S_2},&
t_2&=\overline{t_1}.
\label{eq:singularities}
\end{align}
The residues are computed from the generated matrix.  Their characteristic
polynomials are
\begin{align}
\chi_{t_0}(\lambda)&=\lambda^2(\lambda+1)^2,\nonumber\\
\chi_{t_1}(\lambda)
  =\chi_{t_2}(\lambda)
  &=\lambda(\lambda-1)(\lambda+1)(\lambda+2).
\label{eq:finite-charpoly}
\end{align}
The residue at infinity is minus the finite-residue sum, and
\begin{equation}
\chi_\infty(\lambda)
 =\lambda^4-6\lambda^3+7\lambda^2+6\lambda+(27\sigma-8),
\qquad
\sigma=\frac{S_2}{S_1^2}.
\label{eq:infinity-charpoly}
\end{equation}
These identities are obtained in one generic exact computation over
\(\mathbb{Q}(\sqrt3,i)(m_1,m_2)\).  Specialized rational mass cases are
regressions, not substitutes for the generic derivation.

\section{Positive-mass Galois obstruction}
\label{sec:galois}

Before the normalization \(m_3=1\),
\[
\sigma=
\frac{m_1m_2+m_2m_3+m_3m_1}{(m_1+m_2+m_3)^2}.
\]
Positivity gives \(\sigma>0\), while
\begin{equation}
\frac13-\sigma=
\frac{(m_1-m_2)^2+(m_2-m_3)^2+(m_3-m_1)^2}
     {6(m_1+m_2+m_3)^2}\geq0.
\label{eq:sigma-domain}
\end{equation}
Thus \(0<\sigma\leq1/3\), with equality only at equal masses.

\begin{lemma}[Non-Abelian identity component]
\label{lem:nonabelian}
For every finite positive mass triple, the identity component of the
differential Galois group of the normal variational equation
\eqref{eq:nve-matrix} is non-Abelian.
\end{lemma}

\begin{proof}
The spectra in \eqref{eq:finite-charpoly} are integral, so every finite local
monodromy has all eigenvalues equal to one and is unipotent.  Each unipotent
element of a complex linear algebraic group lies in its identity component.
Assume that identity component is Abelian.  The finite monodromies then
commute, and their product is unipotent.  The fundamental-group relation
would make the monodromy at infinity unipotent as well.

The infinity exponents obtained from \eqref{eq:infinity-charpoly} obey
\begin{equation}
 \big((2\lambda-3)^2-13\big)^2=144(1-3\sigma).
\label{eq:exponent-relation}
\end{equation}
Unipotent infinity monodromy would require all four exponents to be integers.
For two nested-radical roots put \(x=2\lambda_1-3\) and
\(y=2\lambda_2-3\).  Then \(x,y\) are odd integers with
\(x^2+y^2=26\).  The only odd-square decomposition is \(1+25\), which
forces the inner square root in \eqref{eq:exponent-relation} to be \(12\)
and hence \(\sigma=0\).  This contradicts \eqref{eq:sigma-domain}.
Therefore infinity monodromy is not unipotent and the assumed Abelian
identity component is impossible.
\end{proof}

\begin{proposition}[Reproduced planar obstruction]
\label{prop:planar}
For every finite positive mass triple, the declared translation- and
rotation-reduced planar Newtonian Hamiltonian at the fixed nonzero angular
momentum is not completely integrable by meromorphic first integrals in a
complex neighborhood of the parabolic Lagrange phase curve.
\end{proposition}

\begin{proof}
Complete meromorphic Liouville integrability would, by the
Morales--Ramis necessary condition, make the identity component of the
normal variational differential Galois group Abelian
\cite{moralesramis2001}.  Lemma~\ref{lem:nonabelian} gives the contradiction.
\end{proof}

The first-normal-variational polynomial-invariant values
\(\{1/3,8/27,2/9\}\) are not exceptions to this complete-integrability
obstruction.  The positive-mass exceptional set for Proposition
\ref{prop:planar} is empty.

\section{Spatial and KS-atlas transfer}
\label{sec:transfer}

\subsection{Spatial variational subquotient}

Let \(U\) be the collision-free regular part of the twelve-dimensional
translation-reduced spatial phase space.  In pair coordinates its Hamiltonian
has the form
\begin{equation}
H_{\mathrm{sp}}=
\frac{|W_1|^2}{2\mu_1}+\frac{|W_2|^2}{2\mu_2}
+\frac{W_1\cdot W_2}{m_c}
-\frac{k_1}{|X_1|}-\frac{k_2}{|X_2|}
-\frac{k_3}{|X_1-X_2|}.
\label{eq:spatial-H}
\end{equation}
Reflection across the physical plane is a symplectic involution preserving
\eqref{eq:spatial-H}.  Along a planar solution, the in-plane and vertical
variational subspaces are invariant.  The in-plane block is the
translation-reduced planar variational equation.

The planar rotation momentum is
\[
L_z=X_{1x}W_{1y}-X_{1y}W_{1x}
   +X_{2x}W_{2y}-X_{2y}W_{2x}.
\]
Rotational invariance gives \(\{H_{\mathrm{sp}},L_z\}=0\).
At the regular fixed nonzero momentum of the phase curve, the invariant
tangent space \(\ker dL_z\), quotiented by the rotation generator
\(X_{L_z}\), gives the reduced planar variational equation
\cite{marsdenweinstein1974}.  Energy restriction and flow quotient then
give the four-dimensional normal module:
\[
\begin{aligned}
\text{spatial VE }12
&\longrightarrow \text{planar submodule }8
\longrightarrow \text{fixed-}L_z\text{ tangent }7\\
&\longrightarrow \text{rotation quotient }6
\longrightarrow \text{energy tangent }5
\longrightarrow \text{normal quotient }4.
\end{aligned}
\]
The Galois representation of a differential subquotient is an algebraic
image of the full representation.  An Abelian identity component has only
Abelian algebraic images.  Lemma~\ref{lem:nonabelian} therefore makes the
identity component of the spatial variational Galois group non-Abelian.

\subsection{KS cotangent reduction}

For one KS spinor \(q=(q_1,q_2,q_3,q_4)\), use the convention
\cite{kustaanheimostiefel1965}
\begin{align}
\kappa(q)&=\bigl(
  2(q_1q_3+q_2q_4),
  2(q_1q_4-q_2q_3),
  q_1^2+q_2^2-q_3^2-q_4^2
\bigr),\nonumber\\
r&=q\cdot q,\qquad
Y(q)=(-q_2,q_1,-q_4,q_3).
\label{eq:ks-map}
\end{align}
The exact identities
\[
|\kappa(q)|^2=r^2,\quad
D\kappa D\kappa^{T}=4r\Id_3,\quad
D\kappa^{T}D\kappa=4r\Id_4-4YY^{T},\quad
D\kappa\,Y=0
\]
show that the configuration kernel is the circle direction generated by
\(Y\).  The cotangent momentum map and physical projection are
\begin{equation}
\XiMap(q,P)=Y(q)\cdot P,\qquad
x=\kappa(q),\qquad
p=\frac{D\kappa(q)P}{4r}.
\label{eq:cotangent-projection}
\end{equation}
Direct differentiation gives canonical projected brackets on
\(\XiMap=0\); away from that level, the momentum brackets have the expected
magnetic terms proportional to \(\XiMap\).  Applying
\eqref{eq:cotangent-projection} to two spinors gives the regular dimension
chain
\[
18\longrightarrow16\longrightarrow14
\]
for zero-level restriction and quotient by \(\T^2\), including the temporal
canonical pair.

The regularized extended Hamiltonian is
\begin{equation}
\mathcal K=g(H_{\mathrm{sp}}+p_t),\qquad
g=r_{12}r_{23}r_{31},\qquad
\frac{dt}{ds}=g.
\label{eq:sundman}
\end{equation}
For a projectable function \(F\) with no explicit \(t\)-dependence,
\[
\{F,\mathcal K\}
=g\{F,H_{\mathrm{sp}}\}
 +(H_{\mathrm{sp}}+p_t)\{F,g\}.
\]
On \(\mathcal K=0\), with \(g\ne0\), conservation for \(\mathcal K\) is
equivalent to conservation for \(H_{\mathrm{sp}}\).  No off-shell equivalence
is claimed.

\begin{definition}[Physically admissible KS family]
A complete candidate family consists of
\[
(\mathcal K,p_t,\XiMap_1,\XiMap_2,F_1,\ldots,F_5).
\]
It must be pairwise commuting, have rank nine on a relative dense-open subset
of \(\{\mathcal K=\XiMap_1=\XiMap_2=0\}\), be invariant under both circle
actions, agree on qualified chart overlaps, and be independent of the
regularizing clock.
\end{definition}

The shell-relative rank condition is essential.  In adapted canonical
coordinates, a candidate containing \(\XiMap_1p_6\) can have ambient rank
nine away from \(\XiMap_1=0\) while losing its physical direction on the
shell.  The qualification contains this exact negative control.

\begin{lemma}[KS shell transfer]
\label{lem:ks-transfer}
Every physically admissible KS family descends on the regular shell to six
independent commuting physical functions
\((H_{\mathrm{sp}},f_1,\ldots,f_5)\).
\end{lemma}

\begin{proof}
Commutation with \(p_t,\XiMap_1,\XiMap_2\) removes the temporal-angle
components of \(dF_i\).  On the shell write
\[
dF_i=A_i+b_i\,dp_t+c_i^a\,d\XiMap_a,\qquad
d\mathcal K=g(dH_{\mathrm{sp}}+dp_t)+k^a\,d\XiMap_a.
\]
If the descended covectors
\(dH_{\mathrm{sp}},d(A_i-b_iH_{\mathrm{sp}})\) were dependent, the displayed
identities would lift that dependence to one among
\[
d\mathcal K,\ dp_t,\ d\XiMap_1,\ d\XiMap_2,\ dF_1,\ldots,dF_5,
\]
contradicting shell-relative rank nine.  The Poisson product rule from
\eqref{eq:sundman} similarly transfers conservation and involution.  Circle
invariance and overlap equality make the descended functions single-valued
on the physical quotient.
\end{proof}

\begin{theorem}[Qualified spatial and KS obstruction]
\label{thm:main}
For every finite positive mass triple, the complexified
translation-reduced spatial Newtonian three-body Hamiltonian is not
completely integrable by meromorphic first integrals in a neighborhood of
the embedded parabolic Lagrange phase curve at fixed nonzero angular
momentum.
Consequently, the qualified multi-chart KS atlas admits no complete
shell-regular physically admissible meromorphic Liouville family on any
domain containing the regular lift of that neighborhood.
\end{theorem}

\begin{proof}
The normal module in Lemma~\ref{lem:nonabelian} is a differential subquotient
of the spatial variational module.  The Morales--Ramis condition therefore
rules out complete local meromorphic Liouville integrability of the spatial
Hamiltonian.  If a physically admissible KS family existed on a domain
containing the regular lifted neighborhood, Lemma~\ref{lem:ks-transfer}
would descend its restriction to a complete physical family there,
contradicting the spatial obstruction.
\end{proof}

The KS consequence requires shell regularity, zero \(\XiMap\), gauge
projectability, chart compatibility, and clock independence.  Omitting any
one of these conditions changes the model and is outside Theorem
\ref{thm:main}.

\section{Computational qualification}
\label{sec:qualification}

The proof package separates scientific inputs, production certificates, and
a source-disjoint symbolic reconstruction.  Each predecessor is accepted
only by qualification identifier and exact SHA-256 hash.  Status-only
records are rejected.  The principal records are summarized in
Table~\ref{tab:records}.

\begin{table}[htbp]
\centering
\small
\begin{tabular}{@{}lll@{}}
\toprule
Campaign & Role & Qualified status\\
\midrule
WP2-MULTIQ-v1 & Two-chart KS model and transitions & PASS\\
WP4-MR-REPL-v1 & Generic planar NVE reproduction & PASS\\
WP4-LIFT-NI-v3 & Spatial and KS transfer & PASS\\
WP1-CAT-N3BP-v1 & External periodic-orbit catalogue validation & PASS\\
WP1-RESIDUAL-MAP-v1 & Follow-on residual allocation & PASS\\
WP3-MULTI-CAND-v2-KQ-v1 & Frozen-library bounded null search & BOUNDED NULL\\
\bottomrule
\end{tabular}
\caption{Qualified inputs.  The last three rows provide computational
context only and are not premises of Theorem~\ref{thm:main}.}
\label{tab:records}
\end{table}

The generic symbolic replay performs the following decision-bearing checks:
\begin{enumerate}
  \item exact zero-energy and Hamilton-equation tangency of the phase curve;
  \item construction and Hamiltonian identities of the six- and
    four-dimensional variational systems;
  \item symmetry of the Whittaker Hessian;
  \item Fuchsian pole support and simple-pole valuations;
  \item exact finite and infinity characteristic polynomials;
  \item the positive-mass sum-of-squares identity; and
  \item byte-independent equality of the complete normalized result payload.
\end{enumerate}
Elapsed time is excluded from the payload comparison.  Four exact rational
mass cases and a dependency-free production implementation provide
additional convention and regression checks.

Adversarial controls mutate residue polynomials, source hashes, theorem
wording, phase-curve scope, and prerequisite records.  They also attempt to
promote bounded discovery evidence into a theorem premise, accept an
ambient-only KS rank family, accept an internal or anonymous review as
independent, and authorize submission with incomplete metadata.  Every such
mutation must be rejected.

The broader legacy suite is partitioned into deterministic fast tests, slow
simulations, and explicitly quarantined historical prototypes.  That suite
is software-regression evidence only; it does not establish the theorem.

\section{Bounded discovery context}
\label{sec:context}

The WP1 catalogue track validates the project orbit pipeline against an
external unequal-mass periodic-orbit dataset derived from the catalogue of
Li, Li, and Liao \cite{lililiao2021}.  This supports the numerical efficacy
of the orbit tooling under a frozen convention and tolerance contract.  It is
not evidence for the differential-Galois conclusion.

WP3 searched a frozen library of \(285\) candidate functions after quotienting
the known \(L^2\) class.  No promoted invariant survived the locked
qualification.  The WP3 result is bounded to its frozen library and is not
used in the proof.  It does not constitute a search over all functions and
does not support an absence-of-integrals claim.

The two tracks answer different questions.  WP1 tests whether the numerical
pipeline reproduces external orbit evidence.  WP3 allocates a finite
candidate search.  Theorem~\ref{thm:main} instead follows from an exact local
normal-variational obstruction.

\section{Limitations and nonclaims}
\label{sec:limitations}

The theorem concerns complete meromorphic Liouville integrability locally
near one phase curve.  The result does not exclude one additional integral.
It does not exclude non-meromorphic first integrals, integrability on a
domain disjoint from the obstructed neighborhood, or a different notion of
integrability.  It does not cover zero angular momentum, nonzero
\(\XiMap\) magnetic dynamics, gauge-dependent or clock-dependent functions,
or triple-collision regularization.

The computation begins with a published reduced Hamiltonian convention.  It
checks exact equivalence to that convention but does not independently
reconstruct every canonical translation and rotation map from the
unreduced eighteen-dimensional state.  The spatial transfer is an explicit
differential subquotient argument; it is not the invalid inference that
nonintegrability of an invariant subsystem alone implies nonintegrability of
the ambient system.

No global orbit classification follows.  The only domain-level corollary is
by restriction: a complete admissible family on a larger domain containing
the obstructed regular neighborhood would restrict to a forbidden local
family.

\section{Reproducibility}
\label{sec:reproducibility}

The repository contains the manuscript source, bibliography, claim ledger,
proof-dependency graph, canonical artifact manifest, exact predecessor
hashes, and the source-disjoint SymPy reconstruction.  From a clean checkout,
the core package is exercised by
\begin{verbatim}
./scripts/reproduce_wp5_pub_v1.sh --core
\end{verbatim}
The command verifies the specification lock and all pinned inputs, rebuilds
the PDF, reruns the generic symbolic derivation, compares its complete
decision-bearing payload after removing elapsed time, executes the
publication-specific adversarial tests, and emits a machine-readable report.
A full mode additionally runs the partitioned legacy baseline.

Primary papers are identified by bibliographic metadata, URL, and SHA-256
download hash but are not redistributed.  External review is a separate,
fail-closed gate: two independent human reviewers must collectively cover
the planar/differential-Galois and symplectic-reduction/KS arguments, and at
least one clean-room reproduction must occur outside project-controlled
hardware.  No internal process, automated agent, or self-review is treated
as that evidence.

\section{Conclusion}
\label{sec:conclusion}

The package reproduces the published planar Morales--Ramis calculation from
the reduced Hamiltonian and verifies its positive-mass non-Abelian Galois
obstruction with exact arithmetic.  A separately qualified subquotient,
cotangent-reduction, Sundman, and shell-rank argument transfers that
obstruction to the translation-reduced spatial Hamiltonian and to the
declared KS atlas.  This establishes the local, complete-meromorphic
nonintegrability statement in Theorem~\ref{thm:main}, subject to every stated
model boundary.  External mathematical review, independent clean-room
reproduction, author metadata, license choices, and release authorization
remain publication gates rather than inferred facts.

\bibliographystyle{plain}
\bibliography{references}

\end{document}
